\documentclass[10pt,a4paper]{amsart}
\usepackage[margin=19mm]{geometry}
\usepackage{amsmath,amssymb,mathtools}
\usepackage[scr=rsfs,cal=euler]{mathalfa}
\usepackage{xfrac}
\usepackage[unicode,hidelinks,bookmarks=false]{hyperref}
\newcommand{\ie}{i.e.\ }
\newcommand{\cf}{cf.\ }
\newcommand{\resp}{respectively\ }
\newcommand{\Subsection}{\subsection}
\providecommand{\nobreakdash}{\nobreak\hbox{-}}
\setlength{\emergencystretch}{2em}
\multlinegap0pt
\usepackage{bm}
\usepackage{bbm}
\usepackage{dsfont}
\usepackage{xspace}
\usepackage{cancel}
\usepackage{mathtools}
\usepackage{amsthm}
\makeatletter
\@ifundefined{lemma}{\newtheorem{lemma}{Lemma}}{}
\makeatother
\title[Maximal dissipation and shallow water flow -- the dam-break ]{Maximal dissipation and shallow water flow -- the dam-break problem}
\author{Dalal Daw, Marko Nedeljkov, Sanja Ru\v{z}i\v{c}i\'c}
\date{Source: arXiv:2502.08329v2 (CC BY 4.0)}
\begin{document}
\maketitle

\noindent\emph{Source and licence.} Maximal dissipation and shallow water flow -- the dam-break problem. Version arXiv:2502.08329v2 (CC BY 4.0), distributed under the Creative Commons Attribution 4.0 International licence, as recorded in the arXiv licence metadata for this exact version and shown on the abstract page https://arxiv.org/abs/2502.08329v2. Full rights notice: licenses/arxiv-2502.08329-CC-BY-4.0.txt.\par\smallskip

\noindent\textit{Editorial scope.} Counted unit: the Lemma whose source label is \texttt{lemma2} in the article (source0.tex lines 744--765, source label \texttt{lemma2}), delivered together with the article's own complete original proof of it (source0.tex lines 767--902, one \texttt{proof} environment of 136 lines). No mathematics is written, repaired or re-derived: statement and proof are the article's own text line for line. Required context retained uncounted: r2r (lines 705--708). Every definition, display and label of those lines is retained unchanged. Omitted from this package and not redistributed: 1--704, 709--743, 766--766, 903--1334. Nothing inside a retained excerpt is omitted or altered: every retained body is delivered exactly as the article's lines are written, so no omission receipt of any kind accompanies this package. Citations: the retained excerpts carry no citation command at all, so no bibliography entry of the article is transcribed and no reference is redistributed. Reference adaptation: none was needed and none was made; every reference inside the delivered unit resolves inside it, so no unresolved reference or citation is produced. Printed slips kept literal: no printed word, symbol, hypothesis or number of the counted statement or of its proof was altered. Layout and class: the article's own class is \texttt{article} with packages including amsmath, amssymb, amsthm, graphicx, etoolbox, color, verbatim; the wrapper is the lane's pinned \texttt{amsart} editorial wrapper at 10pt on A4, which restates the theorem-like environments the retained text needs and adds the article's own macro definitions that the retained text needs (none), with extra wrapper packages (bm, bbm, dsfont, xspace, cancel, mathtools). The delivered numbering is the unit's own. Assets: no retained span contains a figure environment, a graphics-inclusion command, a PDF-page inclusion, a table or a TikZ picture; no image file is copied into this package, the package declares no asset dependency and no image file is redistributed with it. Licence: version arXiv:2502.08329v2 of this article is distributed under the Creative Commons Attribution 4.0 International licence, as recorded in the arXiv licence metadata for this exact version and shown on the abstract page https://arxiv.org/abs/2502.08329v2. A full rights notice is delivered at licenses/arxiv-2502.08329-CC-BY-4.0.txt. Characters: every retained excerpt is pure ASCII, so the delivered engine, XeLaTeX with its default Latin Modern text font, reports no missing character.
\begin{equation} \label{r2r}
\lambda_{2}{(h}_{0},u_{0}) = u_{0} + \sqrt{gh_{0}} < 0, 
\text{ and  } u_{0} < - \sqrt{gh_{0}}<0.
\end{equation}

\begin{lemma} \label{lemma2}
Suppose that $u_{l}\geq 0$. A solution for $x<0$ may consist of
the following waves and their combinations.
\begin{enumerate}
\item A single $S_{1}$ when
\begin{equation} \label{SR2}
h_{0} > \max \left\{ h_{l}, \; \frac{h_l}{2}\left( \sqrt{1 + 8z} - 1 \right)\right\},
\quad z:=\frac{u_l^2}{gh_l}
\end{equation}
\begin{equation} \label{S}
u_{0} = u_{l} - \sqrt{\frac{g}{2} \left( \frac{1}{h_{0}} + \frac{1}{h_{l}} \right)}
\left( h_{0} - h_{l} \right), \; h_{0} > h_{l}.
\end{equation}
\item A combination $S_{1}+R_{2}$, where $u_{0}$ satisfies
\begin{equation} \label{SR}
u_{l} - \sqrt{\frac{g}{2}\left( \frac{1}{h_{l}} + \frac{1}{h_{0}} \right)}
\left( h_{0} - h_{l} \right) < u_{0} 
< -\sqrt{gh_{0},}\;  h_{0} > h_{l},
\end{equation}
and $h_{0}$ satisfies (\ref{SR2}).
\end{enumerate}
\end{lemma}

\begin{proof}
Let us now check all combinations of elementary waves.

\noindent \textbf{1.}
In the case $R_{1} + R_{2}, $ we have to solve the following   system
\begin{equation*}
\begin{split}
& u_{\ast} = u_{l} + \sqrt{gh_{l}} - \sqrt{gh_{\ast}},\; h_{\ast}< h_{l} \\
& u_{0} = u_{\ast} - \sqrt{gh_{\ast}} + \sqrt{gh_{0}},\; h_{0}>h_{\ast}
\end{split}
\end{equation*}
with respect to the variables \(\ (h_{\ast},u_{\ast})\). Substitution of the first in the
second equation gives
\[u_{0} = u_{l} + \sqrt{gh_{l}} - 2\sqrt{gh_{\ast}} + \sqrt{gh_{0}},
\text{ where } h_{\ast}\leq \min \{h_{l},h_{0}\}\]
and that implies \( u_{0} \leq u_{l} + \sqrt{g}
\left( \sqrt{h_{0}} + \sqrt{h_{l}} \right)\). 
The condition on \({h}_{\ast}\) implies
\[u_{0} \geq u_{l} + \sqrt{gh_{0}} - \sqrt{gh_{l}},\text{ and } 
u_{0} \geq u_{l} + \sqrt{gh_{l}} - \sqrt{gh_{0}}.\]
These two bounds, taken together, may be written as
\[u_{0} - u_{l} > \left| \sqrt{g}\left( \sqrt{h_{0}} 
- \sqrt{h_{l}} \right) \right|.\]
These bounds give \(u_{0} > u_{l}\) that contradicts (\ref{r2r}),
since \(u_{l} \geq 0\). Thus, it is impossible to have a combination of 
two rarefaction waves on the left-hand side.

\noindent \textbf{2.}
To get \(S_{1} + R_{2}\ \ \)solution, one has to find
\((h_{\ast},u_{\ast})\) such that
\[ 
u_{\ast} = u_{l} - \sqrt{\frac{g}{2h_{l}h_{\ast}}\left( h_{\ast} - h_{l} \right)
\left( h_{\ast}^{2} - h_{l}^{2} \right)}, \text{ with }
u_{l} - \sqrt{\frac{g\left( h_{\ast} + h_{l} \right)h_{\ast}}{h_{l}}}<0, \; 
h_{\ast} > h_{l},
\]
and 
\[
u_{0} = u_{\ast} - \sqrt{gh_{\ast}} + \sqrt{gh_{0}}, \text{ with } 
u_{0} + \sqrt{gh_{0}} < 0, \; h_{0} > h_{\ast}.
\]
Again, substituting the first equation in the second one,
\[
\begin{split}
u_{0} = & u_{l} - \sqrt{\frac{g\left( h_{\ast} + h_{l} \right)}{{2h}_{l}h_{\ast}}}\left( h_{\ast} 
- h_{l} \right) + \sqrt{gh_{0}} - \sqrt{gh_{\ast}}\\
= & u_{l} - \sqrt{\frac{g}{2}\left( \frac{1}{h_{l}} + \frac{1}{h_{\ast}} \right)}\left( h_{\ast} 
- h_{l} \right) + \sqrt{gh_{0}} - \sqrt{gh_{\ast}},
\end{split}
\]
where \(h_{l} < h_{\ast} <  h_{0} \). That implies \(h_{0} > h_{l}\). Then, 
assuming that $u_{0}=u_{0}(h_{\ast})$, 
\[
\begin{split}
\frac{{du}_{0}}{{dh}_{\ast}} = & \frac{\frac{g}{2} \frac{1}{h_{\ast}^{2}}}{2\sqrt{\frac{g}{2}
\left( \frac{1}{h_{l}} + \frac{1}{h_{\ast}} \right)}}\left( h_{\ast} - h_{l} \right) 
- \sqrt{\frac{g}{2}\left( \frac{1}{h_{l}} + \frac{1}{h_{\ast}} \right)} 
- \sqrt{\frac{g}{2}.\frac{1}{h_{\ast}}} \\ 
= & \frac{\frac{g}{2}.\frac{1}{h_{\ast}^{2}}}{2\sqrt{\frac{g}{2}
\left( \frac{1}{h_{l}} + \frac{1}{h_{\ast}} \right)}}\left( \frac{1}{h_{\ast}} 
- \frac{h_{l}}{h_{\ast}^{2}} - \frac{1}{h_{l}} - \frac{1}{h_{\ast}} \right) 
- \sqrt{\frac{1}{h_{\ast}}\left( \frac{1}{h_{l}} + \frac{1}{h_{\ast}} \right)} < 0.
\end{split}
\]
So, \(u_{0}\left( h_{\ast} \right)\) is a decreasing function with respect to
\(h_{\ast}\). Then
\(u_{0}(h_{0}) < \ u_{0}(h_{\ast}) 
< u_{0}(h_{l}).\)
The value $u_{0}$ is already bounded from above by $-\sqrt{gh_{0}}<0$.
Writing these two bounds together,
\begin{equation*} 
u_{l} - \sqrt{\frac{g}{2}\left( \frac{1}{h_{l}} + \frac{1}{h_{0}} \right)}
\left( h_{0} - h_{l} \right) < u_{0} 
< -\sqrt{gh_{0},}\;  h_{0} > h_{l}.
\end{equation*}
The fact that shock speed is negative gives another bound on 
\({u}_{0}\)
\[
u_{l} - \sqrt{\frac{g\left( h_{\ast} + h_{l} \right)h_{\ast}}{2h_{l}}}<0 
\text{ that implies }  gh_{\ast}^{2} + gh_{l}h_{\ast} - 2h_{l}u_{l}^{2} > 0.
\]
The roots of that quadratic function are
\[
(h_{\ast})_{1,2} = \frac{- gh_{l} \pm \sqrt{g^{2}h_{l}^{2} + 8gh_{l}u_{l}^{2}}}{2g}.
\]
Thus, 
\[
h_{\ast} > \frac{h_l}{2}\left( \sqrt{1+8z} - 1\right) =: \tilde{h}, \quad \text{where } z:= \frac{u_l^2}{gh_l}.
\]

The value of $h_{\ast}$ is between $h_{l}$ and $h_{0}$, so condition
for its existence is  
\begin{equation*} 
h_{0} > \frac{h_l}{2}\left( \sqrt{1+8z} - 1 \right), \text{ and } h_{0}>h_{l}.
\end{equation*}
The relation $\tilde{h}>h_{l}$ holds for
$z > 1$. If not, the previous bound 
$h_{0}>h_{l}$ is enough for existence of $S_{1}+R_{2}$ solution.

\noindent \textbf{3.}
An $S_{1}$ is much simpler to analyse than 
the previous wave combination. 
The condition for $h_{0}$ is (\ref{SR2}), while $u_{0}$ has to satisfy  
\begin{equation*} 
u_{0} = u_{l} - \sqrt{\frac{g}{2} \left( \frac{1}{h_{0}}+\frac{1}{h_{l}} \right)}
\left( h_{0} - h_{l} \right), \quad h_{0} > h_{l}.
\end{equation*}

\noindent \textbf{4.}
Suppose that there is a single $R_{1}$ solution on the left-hand 
side. That is,  
\(u_{0} = u_{l} + \sqrt{gh_{l}} - \sqrt{gh_{0}}\), \(h_{0}<h_{l}\). 
The condition $\lambda_{1}(h_{0},u_{0})<0$ implies  
\(u_{0} < \sqrt{gh_{0}}\).
From that relations, the condition for \(h_{0}\) given by 
the inequality 
\( u_{l} + \sqrt{gh_{l}} - \sqrt{gh_{0}} < \sqrt{gh_{0}}\). 
It implies \(\sqrt{h_{0}} > \frac{1}{2\sqrt{g}}(u_{l} + \sqrt{gh_{l}})\),
and further on, 
\(h_{0} > \frac{1}{4}\left( \frac{u_{l}}{\sqrt{g}} + \sqrt{h_{l}} \right)^{2} > h_{l}\).
But that contradicts the fact that \(h_{0}<h_{l}\) for $R_1$, so there is no
a single $R_{1}$ solution.

\noindent \textbf{5.}
Finally, suppose that there exists a single $R_{2}$ solution. Then,
\(u_{0} = u_{l} - \sqrt{gh_{l}} + \sqrt{gh_{0}}\), \(h_{0}>h_{l}\),
and \(u_{0} < -\sqrt{gh_{0}}\). The last condition is here because
$\lambda_{2}(h_{0},u_{0})<0$. One could see that $u_{0}>u_{l} \geq 0$ 
since $h_{0}>h_{l}$. But that would mean $u_{0}+\sqrt{gh_{0}}>0$ which
is contradiction with $u_{0}<-\sqrt{gh_{0}}$.

To conclude, the only possible solutions on the left-hand
side of the bed jump (dam) are
\(S_{1}\) with conditions (\ref{S}) and (\ref{SR2}) or 
\(S_{1}+R_{2}\), with satisfied conditions (\ref{SR}) and (\ref{SR2}).
\end{proof}

\end{document}
